\documentclass[12pt,a4paper,onecolumn]{article}
\title{Digest of The New York Times}
\author{The New York Times}
% \definecolor{dark-green}{rgb}{0,0.25,0}
% \usepackage[utf8]{inputenc}

\usepackage{fontspec,xunicode,xltxtra}
\setmainfont[Ligatures=Required,Mapping=tex-text]{Georgia}
% \setmainfont[Ligatures=Required,Mapping=tex-text]{Palatino Linotype}
\setsansfont[Ligatures=Required,Mapping=tex-text]{Myriad Pro}
\setmonofont[Ligatures=Required,Mapping=tex-text]{Courier New}
\usepackage{xeCJK}
\setCJKmainfont{宋体}
\setCJKmonofont{黑体}

% \usepackage{indentfirst}


\usepackage[dvips,dvipsnames,svgnames]{xcolor}
\definecolor{light-gray}{gray}{0.95}

%生成PDF的链接
\usepackage[dvipdfmx]{hyperref}
\hypersetup{
    bookmarks=true,         % show bookmarks bar?
    bookmarksopen=true,
    pdfpagemode=UseNone,    % options: UseNode, UseThumbs, UseOutlines, FullScreen
    pdfstartview=FitH,
    pdfborder=1,
    pdfhighlight=/P,
    pdfauthor={wuxch},
    unicode=true,           % non-Latin characters in Acrobat’s bookmarks
    colorlinks,             % false: boxed links; true: colored links
    linkcolor=blue,         % color of internal links
    citecolor=blue,        % color of links to bibliography
    filecolor=magenta,      % color of file links
    urlcolor=cyan           % color of external links
}
% 使用目录
\usepackage{makeidx}
% 使用图形
\usepackage{graphicx}


\usepackage{titletoc}
\dottedcontents{section}[1.5em]{}{2.3em}{1pc}

% % 页眉页脚的设定
% \usepackage{fancyhdr}
% \pagestyle{fancy}
% \lhead{\includegraphics[height=0.2in]{zte_logo}}

% \usepackage{footmisc}
\usepackage{pfnote}
\usepackage{fnpos}

% % 嵌入的代码显示
% \usepackage{listings}
% \lstset{language=C++, breaklines, extendedchars=false}
% \lstset{keywordstyle=\color{blue},
%         commentstyle=\color{SeaGreen},
%         stringstyle=\ttfamily,
%         showstringspaces=false,
%         backgroundcolor=\color{light-gray}}
% \lstset{basicstyle=\small}

\usepackage{booktabs}

\usepackage{tipa}
\usepackage{phonetic}
\usepackage[gen]{eurosym}

% \usepackage{CJK}
\usepackage{makeidx}
% \usepackage{glossary}

\usepackage{fancyhdr}

\usepackage{lastpage}
\cfoot{\textsf{\thepage\ of \pageref*{LastPage}}}

\pagestyle{fancy}
\fancyhead{}
\renewcommand{\headrulewidth}{0pt}
\pagenumbering{arabic}

\usepackage{graphicx}
\usepackage{picinpar}

\setlength{\parindent}{0in}
\setlength{\parskip}{2.5ex plus4mm minus3mm}
% \setlength{\parskip}{.3\baselineskip}

\setcounter{section}{0}

% bibtex
\usepackage[super]{natbib}
% \usepackage[globalcitecopy,labelstoglobalaux]{bibunits}
% \makeatletter
% \renewcommand{\@bibunitname}{\jobname.\the\@bibunitauxcnt}
% \makeatother

% \renewcommand{\refname}{缩略语}
% \renewcommand{\refname}{}

\makeatletter                   % Reference list option change
\renewcommand\@biblabel[1]{#1:  } % from [1] to 1: 
\makeatother                    %

\newcommand{\citenumfont}[1]{\textsf{#1}}

% \usepackage{tipa}
% \usepackage{phonetic}

\newcommand{\phonetic}[1]{{\fontspec[Colour=CC00CC]{Microsoft Sans Serif} #1}}


%%% Local Variables:
%%% TeX-master: "medical_trivia"
%%% End:

\makeindex
\begin{document}
\date{}

% \maketitle
\tableofcontents
\pagebreak
\section{Unofficial Software Incurs Apple's Wrath}

The iPhone can teach its users how to perform CPR, show them how to mix a White Russian and allow
them to identify any song playing on the radio.

But for some owners of the Apple touch-screen device, the 35,000-plus applications lining the
digital shelves of Apple's App Store are not enough. If you want to use your iPhone as a video
camera, send a photo message or hook it up to your laptop to connect to the Internet, there's no app
for that.

Or at least, no official app.

Through the efforts of developers and hobbyists, the Web is teeming with unauthorized applications
for the iPhone and the iPod Touch (which does everything that the iPhone does except make phone
calls and incur a monthly bill from AT\&T), and there are even some independent online application
stores.

However, in order to use these programs, iPhone owners have to ``jailbreak'' their device --
downloading a bit of software that bypasses Apple's restrictions and allows the installation of
unsanctioned third-party programs.

The growing popularity of jailbreaking has set up a legal battle between Apple, which says it has
the right to regulate what can go on an iPhone, and the users and developers who want to customize
their phones as they see fit.

Jailbreaking is different from unlocking an iPhone, in which users modify the software so the phone
can be used on unauthorized wireless carriers. For some iPhone hobbyists, like Mark Janke,
jailbreaking is akin to customizing a fancy car -- it simply allows owners to personalize the look
of their devices, turning their phones into a brag-worthy accessory and status symbol.

``You can modify your phone and say, 'Hey, look what I did,' '' said Mr.~Janke, who runs a forum
called Hack That Phone, which walks iPhone owners through the jailbreaking process in several
languages, including Swedish and Persian. ``It opens up an amazing world of goodies.''

But according to Apple, jailbreaking is illegal and a breach of the Digital Millennium Copyright
Act. ``These modifications not only violate the warranty, they also cause the iPhone to become
unstable and not work reliably,'' said Natalie Kerris, a spokeswoman for Apple.

In a legal filing with the United States Copyright Office last year, Apple says jailbroken iPhones
rely on modified versions of Apple's operating software that infringe on its copyrights.

In addition, the company says jailbreaking encourages the piracy of approved iPhone applications and
is an expensive burden.``Apple's iPhone support department has received literally millions of
reported incidents of software that crashes on jailbroken iPhones,'' the document says.

Apple filed its brief in response to the Electronic Frontier Foundation's request that the copyright
office recognize an exemption to the digital copyright act that would permit jailbreaking of iPhones
and other devices. The copyright office is expected to rule on the issue by October.

Jailbreaking your own iPhone does not infringe on any copyright, and the tools that help iPhone
owners modify their devices do not distribute anything that belongs to Apple, said Fred von Lohmann,
a senior staff lawyer with the Electronic Frontier Foundation, a nonprofit group that advocates more
openness on the Internet. ``In our view, consumers are allowed to adapt software for their own
personal use,'' he said.

For developers, bypassing Apple's lengthy and opaque approval process allows them to get their
software out quickly and on their own terms. Most iPhone owners who jailbreak their phones do so to
change the interface of the iPhone or add simple features and functions that are not available
through official channels, said Jay Freeman, who operates Cydia, a popular repository for thousands
of third-party iPhone applications and modifications.

For example, Mr.~Freeman's site offers two popular applications that he wrote: Cycorder, a free
program that allows iPhones to record video, and Cyntact, a \$1 program that adds profile pictures
to the iPhone's address book.

Another developer offers a \$9.99 application, iPhoneModem, through Cydia that permits the iPhone to
share its Internet connection with a computer, a practice known as tethering, which cellular
carriers frown upon.

Mr.~Freeman estimates that his platform for apps has been installed on about 2.3 million iPhones and
iPod Touch devices. In April, Apple said it had sold 37 million iPhones and iPod Touches to date.

To counter jailbreaking, Apple releases updates for the operating system software that can render
jailbroken phones useless. But the company's efforts amount to an elaborate game of whack-a-mole as
rogue programmers quickly counter their efforts with their own software updates.

So far, the company has not gone after any specific developers or Web sites that assist people in
jailbreaking their iPhones, said Mr.~von Lohmann, but given Apple's penchant for exerting control
over its products, it could.

The foundation argues that Apple's real goal is to limit competition with its App Store, which has
been wildly successful since Apple opened its digital doors in July. Since then, more than a billion
apps, both free and paid, have been downloaded. Gene Munster, a senior research analyst at the Piper
Jaffray investment firm, estimates that the App Store will have generated roughly \$617 million in
revenue by the end of 2009.

Mario Ciabarra, the operator of a smaller application shop called Rock Your Phone, worries that
Apple could quash competitors by simply including the more popular jailbroken iPhone modifications
in the third version of the iPhone operating system, due out this summer.

Mr.~Ciabarra, whose site offers nine applications for download through a platform he developed
himself, says his products have had half a million downloads since mid-March.

``There's no reason Apple couldn't offer the products we're offering,'' Mr.~Ciabarra said. Like
Apple, Mr.~Freeman and Mr.~Ciabarra take a commission of up to 30 percent from developers selling
software through their stores.

In addition to commissions, Mr.~Freeman says he makes money from selling applications he created,
advertisements on the site, and sponsorships from developers who want their applications featured on
Cydia's home page. All together, he said he earns enough from the site to support himself.

Mr.~Freeman said he was not worried about the competition from Apple or the potential ramifications
of operating an unsanctioned software platform. ``I don't even understand why they're trying to
fight an entire community,'' he said.

William H.~Greene, a professor of economics at New York University who studies digital
entertainment, said most jailbreaking software is free and does not hurt sales of the iPhone. Some
applications available through the independent channels had been rejected by Apple for inclusion in
its store. ``It's hard to see where Apple is being harmed by this,'' he said.

Apple would have more leverage with its copyright claims if it could prove that jailbreaking is
harming its business. ``If the jailbreaker, the individual phone owner, were somehow trying to
profit, then they might have a case,'' Mr.~Greene said. Apple is also well within its rights to
crack down on pirated copies of legitimate applications, he said.

Mr.~Ciabarra, who writes unauthorized apps because the iPhone is ``too fun to pass up playing
with,'' sympathizes with Apple's position -- to a point.

``They're worried about malware being distributed and third-party developers interfering with the
stability of their product,'' he said. ``We just feel users should be able to decide for themselves
if they want to deal with that.''

\section{Obama Seeks Action Against Credit Card Industry}

RIO RANCHO, N.M. -- President Obama stepped up his populist campaign against the credit card
industry on Thursday, pressing Congress to pass new limits on ``anytime, any-reason rate hikes,''
unfair late fees and misleading policies.

As the Senate deliberated on legislation, Mr.~Obama convened a town-hall-style meeting here led off
by a woman who complained of being gouged by her credit card company. The president cited her story
to make the case that banks and credit card companies have been taking advantage of loose rules to
make money at the expense of hard-working Americans.

``We're lured in by ads and mailings that hook us with the promise of low rates while keeping the
right to raise those rates at any time for any reason -- even on old purchases,'' he told hundreds
of people gathered in a high school auditorium. ``You should not have to worry that when you sign up
for a credit card you're signing away all your rights. You shouldn't need a magnifying glass or a
law degree to read the fine print that sometimes doesn't even appear to be written in English.''

Appealing to popular resentment in a time of economic trouble, he cast himself as a reformer taking
on a powerful industry.

``Enough's enough,'' Mr.~Obama said. ``It's time for strong, reliable protections for our consumers.
It's time for reform that's built on transparency and accountability and mutual responsibility,
values fundamental to the new foundation we seek to build for our economy.''

The House has passed legislation imposing new limits on credit card practices, but the Senate has
slowed its consideration of a similar bill. Senate aides said the bill was likely to come to the
floor early next week, and it is expected to pass. The House and Senate would then have to reconcile
their versions before agreeing on a final measure to send to the White House. Mr.~Obama repeated his
demand that Congress send it to him by Memorial Day.

Credit card debt has increased by 25 percent in the last decade, with delinquency rates up more than
a third since 2006, according to statistics cited by the White House. Americans pay \$15 billion in
penalty fees a year, accounting for about 10 percent of the industry's revenues. About one-fifth of
those carrying credit card debt pay more than 20 percent in interest.

The White House arranged for an aggrieved customer who sent the president an e-mail message last
week, Christine Lardner, to introduce him to the crowd here on Thursday. Ms.~Lardner, an Albuquerque
resident, told the crowd that she puts some of her daughters' college costs on a credit card and
when she was approaching the limit, instructed the school to switch payment methods. By mistake, the
school charged another tuition payment to her card, putting her over the limit. The credit card
company then increased her rate to nearly 30 percent from 9.24 percent.

``Raising my rate to 30 percent is ludicrous and corrupt,'' she told the crowd to loud applause. The
policies may be legal, ``but they are ethically wrong and citizens who want to pay their bills are
being taken advantage of.''

\section{Secret Memoir Offers Look Inside China's Politics}

In May 1989, as he feuded with hard-line party rivals over how to handle the students occupying
Tiananmen Square, China's Communist Party chief requested a personal audience with Deng Xiaoping,
the patriarch behind the scenes.

The party chief, Zhao Ziyang, was told to go to Mr.~Deng's home on the afternoon of May 17 for what
he thought would be a private talk. To his dismay, he arrived to find that Mr.~Deng had assembled
several key members of the Politburo, including Mr.~Zhao's bitter foes.

``I realized that things had already taken a bad turn,'' Mr.~Zhao recalls in a secretly recorded
memoir only now coming to light -- a rare first-person account of crisis politics at the highest
levels of the Chinese Communist Party.

From Mr.~Deng's impatient body language and the scathing attacks he received from his rivals,
Mr.~Zhao says in the memoir, which is now being published in book form, it was obvious that Mr.~Deng
had already decided to overrule Mr.~Zhao's proposal for dialogue with the students and impose
martial law.

``It seems my mission in history has already ended,'' Mr.~Zhao recalls telling a party elder later
that day. ``I told myself that no matter what, I would not be the general secretary who mobilized
the military to crack down on students.''

As Mr.~Zhao anticipated, he was immediately sidelined and soon vilified for ``splitting the party.''
He was purged and placed under house arrest until his death in 2005.

But in this long, enforced retirement, it turns out, Mr.~Zhao secretly recorded his own account, on
30 musical cassette tapes that were spirited out of the country by former aides and supporters, of
his rise to national power in the 1980s, his battles with the old guard, and his alliance and
tussles with Mr.~Deng as he loosened Soviet-style controls and helped put China on a path to the
dynamic economic power it has become today.

Mr.~Zhao also tells how he was outmaneuvered during the lengthy student-led pro-democracy
demonstrations in the spring of 1989, setting up his ouster shortly before the military crackdown on
June 4 of that year.

One striking claim in the memoir, scholars who have seen it said, is that Mr.~Zhao presses the case
that he pioneered the opening of China's economy to the world and the initial introduction of market
forces in agriculture and industry -- steps he says were fiercely opposed by hard-liners and not
always fully supported by Mr.~Deng, the paramount leader, who is often credited with championing
market-oriented policies.

In the late 1970s, as the party chief in Sichuan Province, Mr.~Zhao had started dismantling
Maoist-style collective farms. Mr.~Deng, who had just consolidated power after Mao's death, brought
him to Beijing in 1980 as prime minister with a mandate for change. Mr.~Zhao, who like other Chinese
leaders had little training in or experience of market economics, describes his political battles
and missteps as he tried to give more rein to free enterprise.

Roderick MacFarquhar, a China expert at Harvard who wrote an introduction to the new book, said it
had given him a new appreciation of Mr.~Zhao's central role in devising economic strategies,
including some, like promoting foreign trade in coastal provinces, that he had urged on Mr.~Deng,
rather than the other way around.

``Deng Xiaoping was the godfather, but on a day-to-day basis Zhao was the actual architect of the
reforms,'' Mr.~MacFarquhar said in an interview.

Recording over children's songs and Beijing Opera performances on the cassettes in his guarded
compound just north of Tiananmen Square, Mr.~Zhao describes in generally modest terms his tenure as
prime minister and then party secretary.

Mr.~Zhao had initially wrote notes and then around 2000, encouraged by three sympathetic former
officials who were allowed to visit him, decided to tape his memoirs, which he did partly in the
presence of those supporters, said Bao Tong, a former close adviser to Mr.~Zhao who remains under
tight surveillance in Beijing.

Two of the former officials have since died, but one of them, Du Dao-zheng, a former senior official
who oversaw press and publications, arranged for a copy of the tapes to be smuggled to Hong Kong.
Mr.~Du, who lives in China, decided in recent weeks to openly acknowledge his role in a statement
that is quoted in the forthcoming Chinese edition of the memoir but not available in time for the
English edition.

Mr.~Bao, in an interview this week, called the memoir ``very rare historical material'' that
``belongs to all the people of China and to the world.'' He said that the voice was unmistakably
that of Mr.~Zhao and that the memoir's authenticity was not in doubt.

Nearly 20 years after the crackdown and Mr.~Zhao's fall, the edited transcripts are being published
by Simon and Schuster in a book, ``Prisoner of the State: The Secret Journal of Premier Zhao
Ziyang,'' that will be formally released in the United States on May 19. A Chinese-language edition
is being published in Hong Kong.

``This is the first time that such a high Chinese leader has been in a position to tell the truth,''
said Bao Pu, a son of Bao Tong who is an editor of the book and a translator of the English-language
edition. ``At that point, the truth is all he had.''

Also credited as translators and editors are Renee Chiang, a publisher in Hong Kong, and Adi
Ignatius, an American journalist who covered China in the 1980s.

Although the tumult of 1989 is distant for many Chinese, it remains a forbidden subject, heavily
censored on the Internet and rarely if ever mentioned in the state-run media. Beijing authorities
are likely to be unhappy with Mr.~Zhao's airing of inside conflicts as well as his conclusion,
arrived at in isolation after he left power, that China must turn toward parliamentary democracy if
it is to tackle corruption.

In a sharp break with Chinese Communist tradition, even for dismissed officials, Mr.~Zhao provides
personal details of tense party sessions. He attacks several officials, especially his archrival,
the conservative former prime minister Li Peng, who fiercely opposed or, in his view, betrayed him.
He describes how they schemed to turn Mr.~Deng against him.

Mr.~Zhao said that in 1989 he argued that most of the demonstrating students ``were only asking us
to correct our flaws, not attempting to overthrow our political system.''

These efforts to defuse tensions were ``blocked, resisted, and sabotaged by Li Peng and his
associates,'' Mr.~Zhao said.

Perry Link, emeritus professor of Chinese studies at Princeton who was in Beijing in 1989, said:
``Laying bare the personal animosities from such a high position is something new here. It's
certainly the element that will send officials in Beijing through the roof.''

The debate over how to respond to protesting students was part of a continuing struggle over
economic and political change. ``What becomes clear in these tapes is that in the minds of Chinese
leaders, Tiananmen was a continuation of their battles through the 1980s,'' said Bao Pu, who is also
a rights advocate and an editor in Hong Kong.

By forcing out Mr.~Zhao and restoring a political grip that remains largely in place today, the
conservatives squelched hopes that China's economic reforms would be accompanied by systematic
political change. But they were also surprised by the popular revulsion over the crackdown.

With the society in turmoil and especially after seeing the collapse of the Soviet Union, Mr.~Deng
began pressing even harder, in his waning years, for market-style changes, or what he renamed
``socialism with Chinese characteristics.''

Despite his economic triumphs, Mr.~Zhao may be remembered most for his futile effort to head off
violence in 1989. In the tapes, he describes how he learned that the army had started its bloody
march to the square at the heart of Beijing.

``On the night of June 3rd, while sitting in the courtyard with my family, I heard intense
gunfire,'' Mr.~Zhao said.

\section{In Sign of Industry Shift, a Legal Giant Loses 2 Top Partners}

David Fox and Daniel E.~Wolf, two top partners at the New York law firm of Skadden, Arps, Slate,
Meagher \& Flom, have defected to Kirkland \& Ellis in a move likely to send shockwaves through the
Wall Street legal world.

The loss of Mr.~Fox, 51, who was among the highest-paid lawyers at Skadden, is a blow to the firm,
where revenue has fallen across nearly all practice areas. A prominent mergers-and-acquisitions
lawyer, Mr.~Fox is leaving after more than 20 years with the firm, founded in 1948. It is rare for
an established firm to lose such a senior lawyer to a less-known rival, and the move is the first
time a partner in Skadden's New York M.\& A.~practice has jumped to a competitor.

The loss of two noted partners, who together generated tens of millions of dollars in fees annually
for Skadden, could signal a broader shift in the corporate legal landscape as lawyers at large
full-service firms leave for more focused, profitable shops. Last year, Kirkland generated about
\$2.47 million in profit for each partner, a closely watched measurement, compared to \$2.07 million
in profits per partner at Skadden, according to The American Lawyer, an industry magazine.

``Skadden is a terrific firm and has been a home for my entire career,'' Mr.~Fox said in a short
interview. He added that Joseph H.~Flom, the firm's leader, ``is one the greatest lawyers and one of
the greatest people I know.''

Mr.~Wolf, 36, is a rising star in the tight-knit world of deal advisers and, at 30, was one of the
youngest people to make partner at Skadden. He has worked on several prominent mergers with Mr.~Fox,
including BHP Billiton's \$150 billion bid for Rio Tinto last year -- the largest attempted takeover
in history. (BHP ultimately gave up its pursuit of the company.)

While they were partly lured to Kirkland by higher salaries, people close to the private talks said
Mr.~Fox and Mr.~Wolf saw an opportunity to build the firm, based in Chicago, into one of the top
five worldwide M.\& A.~advisers. The firm is currently not in the top 10, according to Thomson
Reuters. Skadden is No.~1.

Kirkland bills itself as younger and more entrepreneurial than some of its long-established
competitors, which also attracted Mr.~Wolf and Mr.~Fox. The firm is also more focused on specific
profit-producing areas like M.\& A.~and bankruptcy, while bigger firms like Skadden, Jones Day, and
Latham \& Watkins offer a vast swath of legal services.

These legal superstores have been cutting costs as business dries up across several practice areas.
This year, Skadden offered to pay its 1,300 associates a third of their base pay not to show up at
work in 2009.

\section{Lessons the Teacher Forgot}

Back in what felt like the golden age of finance, before the fine print of mortgage documents
suddenly became relevant and ordinary people in bars began sharing their worries about credit
default swaps, American banking was celebrated as the envy of the world.

Blue jeans and electronics were arriving from factories scattered from China to Costa Rica, and even
white-collar jobs were slipping overseas, but the sophisticated work of measuring risk and
engineering investments remained the province of the geniuses running Wall Street. Their mastery was
more lucrative than ever, and it was emulated around the globe.

So it registered as a comedown last week to read that Bank of America was selling part of its stake
in the Construction Bank of China, as it scrambled to secure cash in the face of its real
estate-related disasters.

Yes, it has come to this: The largest bank in the United States, putative citadel of free
enterprise, must desperately unload shares in a bank controlled by the Communist Party of China.
That, or risk the wrath of American regulators, newly concerned about how much money financial
institutions have on hand.

Meanwhile, the Treasury last week outlined proposed new rules for derivatives, the exotic
investments whose unsupervised trading was once offered up as a sign of the vibrancy of American
financial innovation, only to become a prime example of how Wall Street set fire to the global
economy.

Not four years ago, when Bank of America paid \$3 billion for a 9 percent stake in Construction Bank
as part of a wave of foreign investment into China, it was supposed to be a sign of Wall Street's
superior money management. American banks -- not just Bank of America, but Citibank, Merrill Lynch
and others -- portrayed their purchases of Chinese institutions as savvy, strategic plays; a way to
get a foothold in the world's largest potential market for seemingly everything.

Still shaking off the cobwebs of its failed experiment in Maoist utopia, China was home to 1.3
billion people whose wallets awaited credit cards, 2.6 billion feet eager for Nike sneakers, and 13
billion fingers waiting to be licked in the thrall of KFC chicken.

American banking executives spoke paternalistically of their Chinese counterparts. Yes, China's
banking system was laced with corruption, but the American banks would bring their culture of modern
finance and teach their new charges how to lend with a dispassionate eye on the bottom line.

``We see value in combining their local knowledge and distribution with our product expertise,
technology and experience with size and scale,'' Bank of America's chief executive, Ken Lewis, said
as he consummated the deal to purchase a piece of Construction Bank in June 2005.

Chinese leaders spoke of their great fortune in gaining Wall Street's tutelage. ``We have much to
learn from our partner in serving customers and creating shareholder value,'' said Construction
Bank's chairman, Guo Shuqing.

These days, of course, talk of Bank of America and shareholder value centers on how much of the
company its newest shareholder -- Uncle Sam -- is destined to own, and whether the bank's shares
retain any value. Bank of America's expertise with size and scale has expanded to encompass the
management of \$45 billion in bailout funds.

For much of Wall Street, the expertise that once was expected to elevate China's financial system
increasingly looks like sorcery, or a vast Ponzi scheme in which banks borrowed vast sums, lent to
virtually anyone, and used incomprehensible models to convince markets that all was fine. They
scattered low-interest credit cards and home equity loan offers like takeout menus, creating the
illusion of prosperity by driving up home values.

In effect, American banks operated not unlike the Chinese banks they were supposed to modernize.
They extracted profits by following a variation of the principle long pursued by their Chinese
counterparts: lend without hesitation while extracting your cut, confident that the government is on
the hook for the losses.

In China, ventures may be spectacularly unprofitable, yet enrich everyone lucky enough to get a
piece. Developers, for example, construct vacant office buildings as an excuse to borrow from state
banks. They rake off a cut for themselves, pay bribes to the party officials who deliver the land
and reward bank functionaries with sumptuous banquets and trips to Macao. Soon enough, the trophy
skyscraper descends into financial disaster, but the developers, bankers and party officials have
already extracted their riches, and for long afterward they will still enjoy them.

Much the same can be said of Countrywide, the mortgage lender that sold itself to Bank of America
last year in a fire sale, after many of its loans went bad. Shareholders were mostly wiped out.
Homeowners suffered foreclosure. But the company's executives made out brilliantly, cashing stock
options amassed during the real estate boom, when Countrywide's share price soared along with its
loan volume. Ditto the Wall Street bankers who enabled Countrywide to lend with abandon by selling
their mortgages to investors.

Now the easy money is gone. Wall Street's financial alchemy has broken down, and bankers are freshly
concerned about the creditworthiness of their borrowers. Bank of America is in such a fix that the
investment it once portrayed as a helping hand to the primitive Chinese banking system must be sold
off in haste just to stay alive.

Shorn of their auras as global paragons of excellence, American banks are even facing pressure to
act more like the Chinese banks they were supposed to reform -- by lending in support of politically
necessary projects.

The biggest criticism of Chinese banks has been that they lend not on the financial merits but in
adherence to the wishes of party leaders. Fearful that a large state company may fail and disgorge
angry, unemployed peasants onto the streets, local party officials pressure state banks to keep the
credit flowing and spare the jobs.

In recent months, the center of the American financial system has effectively shifted from New York
toward Washington, as taxpayer funds keep many institutions in business. Lawmakers and Treasury
officials now implore the banks to use their bailout funds to increase lending, even as the banks
themselves worry about the merits of making loans in a weak economy -- the very conundrum Chinese bankers understand all too well.

Perversely, Bank of America is being forced to shrink its China stake just as it might actually have something to learn about banking from its Chinese partner.

\section{Meetings, Version 2.0, at Microsoft}

This interview of Steven A.~Ballmer, the chief executive of Microsoft, was conducted and condensed
by Adam Bryant.

\textsf{Q. Are there areas you want to improve as a leader?}

A. I race too much. My brain races too much, so even if I've listened to everything somebody said,
unless you show that you've digested it, people don't think they are being well heard. Sometimes you
really don't hear because you're racing. It's just the way my brain works. My brain is just chop,
chop, chop, chop, chop. And so, if you really want to get the best out of people, you have to really
hear them and they have to feel like they've been really heard. So I've got to learn to slow down
and improve in that dimension, both to make me better and to make the people around me better.

\textsf{Q. What's it like to be in a meeting run by Steve Ballmer?}

A. I've changed that, really in the last couple years. The mode of Microsoft meetings used to be:
You come with something we haven't seen in a slide deck or presentation. You deliver the
presentation. You probably take what I will call ``the long and winding road.'' You take the
listener through your path of discovery and exploration, and you arrive at a conclusion.

That's kind of the way I used to like to do it, and the way Bill $[$Gates$]$ used to kind of like to
do it. And it seemed like the best way to do it, because if you went to the conclusion first, you'd
get: ``What about this? Have you thought about this?'' So people naturally tried to tell you all the
things that supported the decision, and then tell you the decision.

I decided that's not what I want to do anymore. I don't think it's productive. I don't think it's
efficient. I get impatient. So most meetings nowadays, you send me the materials and I read them in
advance. And I can come in and say: ``I've got the following four questions. Please don't present
the deck.'' That lets us go, whether they've organized it that way or not, to the recommendation.
And if I have questions about the long and winding road and the data and the supporting evidence, I
can ask them. But it gives us greater focus.

\textsf{Q. How do you assess job candidates?}

A. If they come from inside the business, the best predictor of future success is past success. It's
not 100 percent, but it's a reasonable predictor. For an external candidate, what I've found is that
reference checks are super-important. I didn't used to believe so much in reference checks. You can
always get somebody to say something nice about you. But the truth is, if you ask enough questions
and you ask around, you can really get a profile of who's accomplished various things and who
hasn't.

And I try to figure out sort of a combination of I.Q. and passion. I just ask somebody to tell me
what they've done that they are really proud of and tell me about it. And if it's something you are
proud of, you should be able to answer any question I can come up with, at least at a level that
would satisfy my interest. I ought to be able to see your passion. It might be quiet passion; it
might be bubbly passion. But I should be able to sense that you are one of those people who just
sort of throws themselves into things.

\textsf{Q. Is there a skill or qualification or trait that you're looking for in prospective hires that}
didn't matter as much 10 years ago?

A. Mostly, I'm still looking for what I've always looked for: extremely smart and talented people
who love to work hard, who are passionate about technology and who have a great foundation in math
and science. But compared to 10 years ago, technology is more complex, products and services span
people's lives in new ways, and our business is much more global. So it's more important that people
can think outside the confines of their individual expertise and their product group and connect the
dots between technologies, customer needs and markets in new ways.

\textsf{Q. What's the most challenging part of your job?}

A. Finding the right balance between optimism and realism. I'm an optimist by nature, and I start
from the belief that you can always succeed if you have the right amount of focus combined with the
right amount of hard work. So I can get frustrated when progress runs up against issues that should
have been anticipated or that simply couldn't have been foreseen. A realist knows that a certain
amount of that is inevitable, but the optimist in me always struggles when progress doesn't match my
expectations.

\textsf{Q. Fill in the blank. You want the culture of your company to be more \_\_\_\_\_?}

A. Efficient. The right word is efficient. That's the direction that every business leader is
steering their company culture toward right now. Given the current economic climate and the
uncertainty about how long the recession will last, this is a time when organizations need to do
more with less, and Microsoft is no exception. We've made good progress, but for a company that has
grown every year for more than 30 years, learning how to operate under more constrained
circumstances is not always that easy.

At the same time, the need to be more efficient drives us all toward sharper focus on what is
important and what can truly move the needle in terms of meeting customer needs and taking market
share. Of course, we need to be innovative, but we also need to be efficient.

\textsf{Q. Any books on management and leadership that you've found particularly useful?}

A. Jim Collins's book ``Built to Last.''

\textsf{Q. In all the speeches you've given, is there a favorite line or story or passage or quotation?}

A. In February I was invited to share my business perspective on the economic downturn with House
Democrats at their annual retreat. In that speech, I got to share something that my dad always told
me growing up, which is a simple piece of advice that really shaped my approach to life and to
business.

My dad worked for Ford for 30 years. When I was a kid, he'd say: ``If you're going to do a job, do a
job. If you're not going to do a job, don't do a job.'' What he meant was, if you really want to
accomplish anything, you have to be committed, motivated, tenacious and smart about what you do.
That's really just the essence of the American work ethic, but it's one of the most important things
I ever learned.

\textsf{Q. If you had to choose another profession, what would it be?}

A. Education, probably. I like working with young people, and I think it's really important to
encourage talent. I love basketball, so I could see myself as a high school basketball coach. I
think a basketball team that I coached would have a really good chance of being a winner.

\textsf{Q. What would you like business schools to focus on more, or less?}

A. I'd like to see more emphasis on the importance of taking the long view. Companies focus too much
on short-term results in business. It takes patience to build a great business, and sometimes you
have to be willing to make the long-term investment and then keep at it if you want to succeed.

\textsf{Q. If you could teach any b-school course, including one that you create, what would it be?}

A. Leadership. Microsoft has grown from 30 people to more than 90,000 since I started, so I've had
the chance to play a leadership role at practically every stage imaginable in a company's growth and
development. I've learned a lot about leadership along the way from some great people that I've
worked with and through experience.

I've come to believe that to be a great leader, you have to combine thought leadership, business
leadership and great people management. I think most people tend to focus more on one of those
three. I used to think it was all about thought leadership. Some people think it's all about your
ability to manage people. But the truth is, great leaders have to have a mix of those things.

\section{Utah Governor Chosen as Ambassador to China}

WASHINGTON -- President Obama on Saturday selected one of the nation's leading Republican governors
to serve as the ambassador to China, nominating Jon Huntsman Jr.~of Utah for a diplomatic post that
Mr.~Obama called ``as important as any in the world.''

The decision renewed Mr.~Obama's sometimes rocky efforts to reach across the aisle in an era of
political polarization while signaling that he would continue the American policy of economic and
strategic engagement with China advanced by presidents of both parties in recent times.

``We will launch a new era of partnership between our two nations that will advance our shared
dreams of opportunity and security in America, in Asia and around the world,'' Mr.~Obama said. ``I
can think of no more important assignment than creating the kinds of bridges between our two
countries that will determine the well being not just of Americans and Chinese, but also the future
of the world.''

Governor Huntsman stood at the president's side in the Diplomatic Room of the White House on
Saturday morning as he accepted the nomination, which he said was unexpected. As a co-chairman of
Senator John McCain's presidential campaign last year, Mr.~Huntsman said he did not anticipate being
``called into action by the person who beat us.''

Mr.~Huntsman, 49, learned to speak Mandarin Chinese from his time as a Mormon missionary in Taiwan.
He has worked in the two Bush administrations, serving as ambassador to Singapore in the final year
of President George Bush's term and as a deputy trade ambassador for President George W.~Bush.
Mr.~Huntsman's name was floated as a possible ambassador to China for the second President Bush, who
instead sent Clark T.~Randt Jr., a longtime friend and a lawyer with Asia experience.

Political and diplomatic veterans called Mr.~Huntsman a good choice.

``China is a place where we really need to have more of a bipartisan approach and I think he would
contribute to that,'' said James Sasser, a former Democratic senator from Tennessee who was
ambassador to China under President Bill Clinton. ``I think it does take the politics out of it to a
considerable degree.''

Michael Green, who was the top Asia adviser to the younger Mr.~Bush, said Mr.~Huntsman was well
suited for the challenges ahead. ``U.S.-China relations can get rough and tumble over the next few
years,'' Mr.~Green said, citing Beijing's economic and military expansion. ``All of that is going to
be harder, not easier, so putting someone in who knows China but is seen as firm is a good thing --
particularly someone who can bring along moderate Republicans.''

For Mr.~Obama, the selection of Mr.~Huntsman is something of a political coup. Mr.~Huntsman has
emerged as one of the nation's most visible Republican governors since he won a second term last
year and was expected to at least consider seeking his party's presidential nomination to run
against Mr.~Obama.

As the Republican Party forges through a period of reinvention, Mr.~Huntsman has been a leading
voice for moderation on issues like immigration, gay rights and the environment, even though he has
led one of the most conservative states in the country. While it is uncertain whether he would have
actually sought the Republican nomination, if he is confirmed by the Senate he will presumably be
out of the mix in 2012.

``When the president of the United States asks you to step up and serve in a capacity like this,
that to me is the end of the conversation and the beginning of the obligation to rise to the
challenge,'' said Mr.~Huntsman, who was joined by his wife, Mary Kaye, and the couple's seven
children, one of whom was adopted from China. An adviser said Mr.~Huntsman did not intend to change
his party affiliation.

Rahm Emanuel, the White House chief of staff, oversaw the courtship. He said the White House
believed that the governor was the ``best-qualified person for the job,'' considering his experience
in the region, his fluency in the language and culture, and his knowledge of issues affecting the
region. Mr.~Huntsman's work as ambassador, as well as in the trade representative's office, and his
record as governor impressed Mr.~Obama, Mr.~Emanuel said.

The nomination of Mr.~Huntsman is the latest in a series of Republican appointments in the Obama
administration. The president selected two Republicans to serve in his cabinet, but aides said he
was frustrated when a third, Senator Judd Gregg of New Hampshire, withdrew his nomination as
commerce secretary.

Jeff Bader, the president's senior Asia adviser, who worked in the trade representative's office
with Mr.~Huntsman, suggested the governor for the post.

Mr.~Emanuel said he had three later phone conversations with Mr.~Huntsman, and after the governor
indicated he was interested, the president called him on May 5 to offer him the job. The two met in
the Oval Office on May 9 -- their first substantive face-to-face meeting -- when Mr.~Huntsman was in
town for the White House Correspondents' Association dinner.

% \pagebreak
% \printindex
\section{New Starbucks Ads Seek to Recruit Online Fans}

In a new ad campaign, Starbucks wants to tell its message to a new generation of coffee drinkers and
then recruit them to retell the story online.

The coffeehouse chain is putting up new advertising posters in six major cities. To further spread
its message, it is trying to harness the power of online social networking sites by challenging
people to hunt for the posters on Tuesday and be the first to post a photo of one using Twitter.

The outdoor ads boil Starbucks' message down to headlines, some of them veiled jabs at competitors:
``If your coffee isn't perfect, we'll make it over. If it's still not perfect, you must not be in a
Starbucks.''

The marketing campaign, which started this month in newspapers and magazines, is described by the
company, which generally avoids traditional advertising, as the biggest marketing effort it has
undertaken.

The Starbucks campaign goes up against a major advertising blitz by McDonald's promoting its new
line of McCaf\'e coffee drinks. McDonald's is reportedly spending more than \$100 million on
television, print, radio, billboard and Web ads that show people's bland work days full of commutes
and cubicles being improved with a whipped cream-covered McCaf\'e drink. The promotion is the
biggest product introduction for McDonald's since it started serving breakfast in the 1970s, the
company said.

Starbucks, which would not say how much its campaign will cost, says it is not losing coffee
drinkers to McDonald's. Rather, its customers, in trying to save money, have been cutting back on
its espresso concoctions. Same-store sales were down 8 percent in the first three months of the
year.

The competition, said Terry Davenport, chief marketing officer at Starbucks, ``is trying to just
commoditize coffee and take it down to a level where all coffee's the same, and if coffee's coffee,
you might as well buy the cheap stuff.''

He added, ``We just don't believe that to be true. That's why we wanted to tell our stories.''

Some coffee drinkers in the new generation see Starbucks coffee as a commodity, too, having grown up
with a store on every corner. ``All they know is Starbucks, the big company,'' so the ads seek to
highlight the quality of the coffee, Mr.~Davenport said.

Starbucks's text-heavy ads have bold headlines written on a background that looks like a burlap
coffee sack, meant to evoke roasted coffee, said David Lubars, chief creative officer of BBDO North
America, the agency that created the campaign and part of the Omnicom Group.

The full-page newspaper ads go to some length to describe how Starbucks selects only the best 3
percent of beans and roasts them until they pop twice, and gives its part-time workers health
insurance.

Starbucks chose the copy-filled ads, which were popular in the 1960s and 1970s, because it wanted to
put its full story out, Mr.~Davenport said. ``Even if you cruise by and don't stop to read every
word, the net impression is, 'Wow, Starbucks has a lot to say about coffee.'''

That might not be the right strategy for young people, said Richard Honack, a professor of marketing
at the Kellogg School of Management at Northwestern. Unlike Starbucks' older customers, ``Generation
Y goes to Starbucks for the Internet, the music, a place to hang out,'' he said. ``Selling them the
coffee and where the coffee comes from? I just don't know if that's a good idea.''

The idea for the Starbucks photo contest came from watching what people already do on Facebook and
Twitter, said Chris Bruzzo, vice president for brand, content and online at Starbucks. Each year,
people race to post the first photos of Starbucks shops decorated in red for the holidays, he said,
and on Flickr, people vie to post photos that include multiple Starbucks stores in the same shot.

``It shows a level of connection to our brand that we wouldn't have concocted on our own,''
Mr.~Bruzzo said.

Starbucks has other social media initiatives planned for this campaign, including a contest for
Starbucks store employees to submit headlines for future ads and YouTube videos with coffee experts
talking about Starbucks coffee.

Starbucks says it thinks its campaign will be helped by its 1.5 million fans on Facebook and 183,000
followers on Twitter. On the Saturday before the presidential election, Starbucks sponsored a single
60-second television commercial on ``Saturday Night Live'' advertising a coffee giveaway on Election
Day. Starbucks then posted the video online. By Tuesday, it was the fourth-most-viewed video on
YouTube, and people were mentioning Starbucks on Twitter every eight seconds.

Still, it is difficult to measure the effects of social media -- a follower on Twitter does not
necessarily translate to a daily Frappuccino drinker.

Mr.~Bruzzo said Starbucks' social media presence gave it an advantage over competitors with gigantic
ad budgets because its fans wanted to talk about it online. ``It's the difference between launching
with many millions of dollars versus millions of fans.''

\section{Sri Lankan Rebel Leader Also Served as a Cult Figure}

BANGKOK -- For a quarter century, Vellupillai Prabhakaran led a brutal and committed insurgency that
terrorized Sri Lanka with massacres, suicide bombings and assassinations.

The founder of the Liberation Tigers of Tamil Eelam, Mr.~Prabhakaran, 54, was apparently killed
Monday in a government military offensive that destroyed the remnants of an army he had built that
at one point numbered some 10,000 fighters.

Many analysts have predicted that the movement, with its military destroyed and its territory gone,
will reconstitute itself using the tactics it was built around: guerrilla warfare and terrorist
bombings.

But Mr.~Prabhakaran was as much a cult figure to Tamil separatists as he was a commander, and it was
not clear what would become of his insurgency without him.

Though the Tamils are largely Hindu and although he sometimes employed Marxist terminology,
Mr.~Prabhakaran's movement was for the most part neither religious nor ideological. It focused on
the creation of a Tamil homeland in the area known as Eelam, in northern and eastern Sri Lanka.

For years, he commanded what amounted to a shadow state with its own flag, police and court system.
He built a fanatical army, including many child recruits, and ordered them to abstain from sex and
cut themselves off from personal ties.

Sri Lanka's Tamils, who make up just 12 percent of a population that is largely Buddhist and
Sinhalese, have been marginalized since independence from Britain in 1948.

The dominance of the Sinhalese was cemented with the declaration of Sinhala as the national language
and Buddhism as the principal religion, and with privileges for Sinhalese in education and
government jobs.

Mr.~Prabhakaran was a pioneer in the tactic of suicide bombings, creating a squad called the Black
Tigers that carried out scores of killings over the years.

The bombers attacked both government figures and moderate Tamils and claimed a long list of
high-profile victims, including India's prime minister, Rajiv Gandhi, in 1991 and Sri Lanka's
president, Ranasinghe Premadasa, in 1993.

Mr.~Gandhi was apparently chosen for assassination in revenge for India's military intervention
against the Tigers in Sri Lanka in the late 1980s. India withdrew its troops in May 1990 after 1,200
soldiers were killed fighting the Tigers.

A chubby man with a bushy mustache, Mr.~Prabhakaran (pronounced PRAH-bah-ka-ran) was a master of
concealment and escape, rarely appearing in public.

It was his refusal to compromise that undermined repeated attempts at peace agreements and led to
the final battles that destroyed the movement he had built.

As recently as last November, in his annual public address, he said, ``No sane voice is being raised
either to abandon war or to seek a peaceful solution to the conflict.''

Born to a middle-class family in the Tamil heartland of Jaffna on Nov. 26, 1954, the youngest of
four children, Mr.~Prabhakaran dropped out of school and devoted his life to guerrilla war, going
underground in 1972.

He announced himself with the assassination in 1975 of a governor of Jaffna, walking up to him at a
Hindu temple and shooting him at point-blank range.

He was fascinated by shootouts in Westerns, according to his biographer, the Indian journalist M.R.
Narayan Swamy. ``He would take slow steps with a revolver stuck into his shirt, make a sudden turn,
whip out the revolver and fire at an imaginary enemy,'' Mr.~Swamy quoted a friend as saying. ``He
never got tired of it.''

\section{Ford Rejects Big Cuts in Dealer Network}

DEARBORN, Mich. -- The Ford Motor Company says it will not match big dealer cuts made by General
Motors and Chrysler, and expects to benefit from the elimination of nearly 2,000 competing
showrooms.

The company's director of North American sales, James D.~Farley, said in an interview Monday that
Ford had been pushing steadily to consolidate its dealers, rather than trying to end contracts or
let them expire. Ford has reduced its dealer network by 700 since 2005, leaving it with 3,700
nationwide. Mr.~Farley declined to specify Ford's targets for additional reductions, but said they
would be small compared with the cuts by G.M. and Chrysler. Last week, Chrysler, which is in
bankruptcy, said it would terminate nearly 800 dealers, or 25 percent of its total, by next month.

G.M. said it would cut 1,100 dealers next year in the first phase of a plan to reduce its dealers by
40 percent, to about 3,600 from 5,900.

Ford is not accepting any government bailout money, in contrast to its two crosstown rivals. It is
taking every opportunity to highlight the difference, even as it searches for ways to match some of
the cost-cutting structural changes at G.M. and Chrysler.

Mr.~Farley was particularly critical of the Chrysler plan, noting how it would affect millions of
consumers with little warning.

``It seems very abrupt and unplanned,'' he said. ``You don't orphan four million customers overnight
without some fallout.''

Some of those customers, primarily those in rural areas, will migrate to Ford dealerships, he said.
``It really depends on how G.M. and Chrysler handle these orphan owners,'' he said. ``If they don't
give them a lot of attention, it will result in consumers going to other brands.''

Chrysler has asked its bankruptcy judge to approve the immediate closing of its dealerships at a
hearing in early June.

The company said it would try to help affected dealers sell their inventories, an estimated 44,000
vehicles, to surviving Chrysler dealers.

Mr.~Farley said the way that was handled could affect the overall market for new cars in the short
term.

``We are very concerned how they are going to handle those 44,000 units,'' he said. ``It's like a
liquidation sale now, and the biggest issue is whether they will cut prices to move the inventory.''

He said Ford would not offer greater sales incentives to match any price reductions that may result
from the sell-off of inventories at G.M. and Chrysler. G.M. and Chrysler have said that by reducing
the number of dealerships, the surviving showrooms will sell more vehicles, make more money and be
able to invest more in dealerships and service.

``We don't think it's productive to just get rid of rural dealers,'' Mr.~Farley said. The cuts at G.M. and Chrysler were partly focused on smaller and rural dealers. Chrysler dealers are expected to protest the closings in bankruptcy court. G.M. has not completed plans for dealer reductions. It said last week that it was simply informing 1,100 of its dealers that their franchise agreements would not be renewed next year, and that further decisions on the timing were several weeks away.

By then, G.M. could be following Chrysler into bankruptcy. The company is under a June 1 deadline set by President Obama to restructure its operations and reduce its debt, or be forced into a bankruptcy filing.

\section{Passengers' Advocates See Progress}

After a nearly a decade of failed efforts to win rights for airline passengers, advocates may
finally score some modest gains.

Congress is considering three initiatives: requiring the airlines to provide passengers stuck on
grounded aircraft with food, water, functioning restrooms and medical treatment (and possibly the
option to get off the plane after three hours); establishing a consumer complaints hot line and
publishing the phone number on boarding passes; and requiring airlines to disclose a flight's
on-time record at the point of purchase.

Since the proposals are either part of the Federal Aviation Administration appropriation bills
moving through Congress or are expected to be offered as amendments to that legislation, they have a
better shot at success this time around, proponents say. A vote in the House is expected sometime
this week.

The thorniest issue is whether Congress will impose a time limit on keeping passengers on planes
stuck on the tarmac.

The carriers strongly oppose that prospect, arguing that the logistics involved would make it
difficult to give people the option of getting off. Yet four Canadian airlines recently agreed to
let passengers disembark if a ground delay exceeds 90 minutes -- if it is safe and practical to do
so.

``This is subject to the aircraft commander's discretion,'' said George Petsikas, president of the
National Airlines Council of Canada. ``But clearly, we understand that when things get long on board
we have to look after our passengers' well-being and comfort.''

Mr.~Petsikas acknowledged that the airlines agreed to the new rules, which begin on June 9, partly
as a response to a more aggressive proposal before Parliament that would require carriers to pay
passengers compensation for delays.

While such financial penalties are not under consideration in Congress, the idea is supported by the
National Business Travel Association.

``We think the F.A.A. should take a look at whether refunds could be appropriate for a multitude of
reasons, including delays,'' said Stewart Verdery, the group's government relations consultant.

Mr.~Verdery said the travel association deferred to the discretion of the airlines on how to handle
long onboard delays, as did the Airports Council International-North America.

``From the airport's perspective, we support the airlines making that decision,'' said Deborah
McElroy, an executive vice president with airports council. But she said that airports typically did
have the resources to help get passengers off aircraft, even when a gate was not available, using
buses or stairs or an alternate location like a hangar.

``If the decision is made to deplane passengers, we have plans in place to make that happen,''
Ms.~McElroy said.

One issue that has clouded the debate over how to handle extended tarmac delays is determining how
often these situations occur.

For years, the Department of Transportation primarily tracked flights that experienced long taxi-out
times from the gate to the runway. Using this criteria, about 7,150 flights spent two hours or more
on the tarmac before taking off in 2008, and 1,231 spent three hours or more awaiting take-off.

Beginning last October, the Transportation Department also started tracking data about additional
flights that spent a long time on the ground: flights that were diverted to an alternate airport,
planes that left the gate and returned multiple times, flights that were delayed after landing or
flights that were delayed and then canceled. (International flights and flights on many smaller
regional carriers are still not counted.)

The new parameters have added about 25 flights a month to total delays of three hours or more, but a
group of passenger rights advocates questions whether the Transportation Department is getting
accurate data from the airlines.

According to Kate Hanni, executive director of FlyersRights.org, which has been pushing for better
accounting, a dozen airlines filed data with the department suggesting that they always gave
passengers the option of getting off a plane when a flight was diverted to an unscheduled airport --
a virtual impossibility since records showed many of those flights spent less than half an hour on
the ground.

``All this information comes from the airlines, and unless someone proves that it's wrong, the
D.O.T. reports what they get from the airlines,'' Ms.~Hanni said. ``It calls all of the data into
question.''

Dave Smallen, a spokesman for department's Bureau of Transportation Statistics, said the agency was
planning to meet with Ms.~Hanni and Mark Mogel, a FlyersRights volunteer and software engineer who
has been analyzing the data, to discuss the issues they raised.

David Castelveter, a spokesman for the Air Transport Association, said the airline trade group
believed the data ``accurately captures the extent and cause of delays,'' and continued to oppose
rules about deplaning passengers.

``For safety reasons, you can't commit with absolute authority that you can get somebody off an
airplane,'' Mr.~Castelveter said.

But he said some airlines were already complying with another proposal being considered, which would
require carriers to disclose a flight's on-time record at the point of purchase. On United's Web
site, for instance, customers shopping for a ticket can view ``flight info'' to see a flight's
on-time record.

While that requirement is not currently part of the F.A.A. reauthorization bills, Senator John
Thune, Republican of South Dakota, introduced a similar measure in 2007 (the Informed Air Traveler
Act) and a spokesman said he planned to offer the language as an amendment to the F.A.A. bill.

A provision to establish a consumer complaint hot line, and publish it on boarding passes, is also
being considered. While the Transportation Department's Aviation Consumer Protection Division has a
number passengers can call to report problems like mishandled luggage, it is not widely publicized.

``It would be a monumental step forward to get any of this legislation passed,'' Ms.~Hanni said.

Even so, these measures do not address many of the frustrations customers experience while flying,
or the underlying problem of systemic delays.

``I'm a capitalist,'' said William McCurry, a consultant based in Princeton, N.J., explaining that
he generally supports free-market principles, but joined FlyersRights because of his frustrations
with air travel -- and the inability to rely on the system to meet clients.

``This organization should not need to exist,'' Mr.~McCurry said. ``That's why we have Congress and
the F.A.A. But nobody is looking out for passengers.''

\section{In U.S., Steps Toward Industrial Policy in Autos}

President Obama has cast himself as a reluctant interventionist in two of the nation's major
industries, Wall Street and Detroit. The federal aid, he says, is a financial bridge to a postcrisis
future and the hand-holding will be temporary.

Even so, the scale of the government investment and control -- especially by the auto task force now
vetting plans at Chrysler and General Motors -- points to an approach that has been shunned by the
United States more than other developed nations.

``By any coherent definition, this is industrial policy,'' said Marcus Noland, a senior fellow at
the Peterson Institute for International Economics.

Industrial policy refers to government programs tailored for a specific industry instead of actions
whose effects are felt across an economy, like monetary policy or tax rates.

Mr.~Obama's declaration on Tuesday of tougher rules on automobile emissions and mileage standards is
both an environmental and an industrial policy. The new national standards require carmakers to
produce fleets that are 40 percent cleaner and more fuel efficient by 2016. In the United States,
industrial policy has long been viewed with suspicion by many policy makers and economists, who
consider it government meddling in the private sector and a violation of free-market principles.

Industrial policy was a hotly debated topic in the 1980s, with the rising challenge from ``Japan
Inc.'' in industries as varied as automobiles and semiconductors. But as the competitiveness of the
American economy revived and the Japanese challenge ebbed, the attention to government support
faded.

Today, economists, along with public policy and Japan experts, say the model sometimes has a role,
and they point to qualified success stories. But with Washington venturing into this area, they warn
that industrial policy tends to be strewn with pitfalls, as political influence too often trumps
economic efficiency.

The strategy, they say, works best when the path ahead is well defined: catching up in an industry,
as Japan did so well with semiconductors for a while; or softening the social impact of industries
that are consolidating or contracting, as Japan did in shipbuilding and parts of steel making.

Industrial policy, they say, frequently falters when it approaches a technological frontier.

Japanese industrial policy in computing technology shows both the promise and the limits of the
approach. Japan's government-guided program to develop improved processes for making computer chips
in the late 1970s proved quite successful.

By the early 1980s, the American leader, Intel, was staggering under the Japanese assault so badly
that I.B.M. stepped in and bought a 20 percent stake in Intel to prop up a crucial domestic
supplier. Under attack, Intel made the leap to a new generation of chip technology.

The Japanese chip makers did not make that bold bet on the future. And the country's computer
industry was single-mindedly focused on following I.B.M. in bulky mainframe computers. ``Japan
missed the big shift to microprocessors and personal computers in the 1980s,'' said Edward
J.~Lincoln, an economist and Japan expert at the Stern School of Business at New York University.

Detroit faces two challenges: first shrinking in size and trimming costs, and then moving toward a
future that is expected to rely on alternative technologies, like battery-powered electric cars.

In contracting businesses, industrial policy suggests an active role for government in softening the
effects for displaced workers and affected communities. ``To me, the strongest argument for
industrial policy was always easing the pain and smoothing the adjustment,'' said Robert B.~Reich, a
professor at the University of California, Berkeley.

Mr.~Reich, a labor secretary in the Clinton administration, is critical of the Obama
administration's auto rescue plan for being focused mainly on the financial engineering of recovery.
He recommends wage insurance and retraining programs linked to economic development projects to
bring new industries to communities suffering from auto plant closings.

In the long term, a different and perhaps bright future beckons for Detroit. It could depend on
nonfood biofuels, like wood chips and algae, which experts predict could someday be produced for \$1
a gallon and substitute for up to half of the gasoline consumed by American cars.

Another potentially promising path, energy experts say, is electric cars, using advanced batteries.
Japan leads in plug-in hybrids, and China and other foreign producers have ambitious plans. But
G.M.'s entry, the Chevrolet Volt, is scheduled to go on sale next year. And Ford has a prototype
car, called Project M.

A lithium-ion battery start-up, A123 Systems, has applied for a \$1.8 billion loan from the Energy
Department to build a series of car-battery plants, the first planned for southeast Michigan.

Such a future, experts say, might also entail auto companies transforming the way they do business,
working closely with electric utilities or start-ups like Better Place, based in Palo Alto, Calif.,
to provide cars, batteries and electricity as a bundled service with monthly fees.

``The technology is beginning to appear to make us think very differently about cars,'' said David
E.~Cole, chairman of the Center for Automotive Research, a nonprofit organization in Ann Arbor,
Mich. ``And whole new business models are becoming possible and likely.''

Industrial policy, economists say, can hasten technological shifts by financing research in a field.
Such research, they say, benefits the economy as a whole -- a payoff beyond the spending that
individual companies could justify.

``That is the big argument for public-funded but private sector-conducted research and
development,'' said Robert N.~Stavins, an economist at Harvard. ``You want companies and researchers
finding the way forward instead of having it defined by government.''

But any alternative future is years down the road for Detroit. It will depend on federal research
incentives, environmental policies and fuel prices.

Still, government policy, some experts say, that is merely a lifeline for Detroit is too narrow a
goal. That was what was done at the start of the 1980s, when Washington negotiated so-called
voluntary restraints on auto exports with Japan. The pact eased the competitive threat on Detroit
and forced Japanese automakers to build factories in the United States. Yet it proved a short-lived
reprieve because Detroit's underlying problems -- high costs and dependence on fuel-hungry models --
were not addressed.

``If all we do is put money into Detroit to give them breathing room, it will fail and we shouldn't
do it,'' said Clyde V.~Prestowitz, president of the Economic Strategy Institute and a trade
negotiator in the Reagan administration.

Yet a more comprehensive, industrial-policylike approach to Detroit carries its own perils,
economists say. In trying to manage the industrial shrinkage, they say, there is a fine line between
easing the social impact and protecting jobs in ways that inhibit economic change and renewal. In
pursuit of new growth, governments risk encouraging overinvestment in areas that prove to be
technological dead ends.

In the Japanese experience, economists see evidence of both dangers. Problems, they say, are
typically byproducts of what economists call ``political capture.'' That is, an industrial sector
earmarked for special government attention builds up its own political constituency, lobbyists and
government bureaucrats to serve that industry. They slow the pace of change, and an economy becomes
less nimble and efficient as a result.

Economists say the phenomenon is scarcely confined to nations with explicit industrial policies and
cite the history of agricultural subsidies in America or military procurement practices.

But going down the path of industrial policy certainly holds that risk. ``You have to bear in mind
the opportunity costs of these kinds of government interventions, and remember that life is not an
economic textbook and that politics can easily override economic rationality,'' said Mr.~Noland, an
author, with Howard Pack, of ``Industrial Policy in an Era of Globalization: Lessons From Asia.''

\end{document}
